% !TEX root = ../../../main.tex
\section{Computing eigenvalues and eigenspaces}
\label{sec:linear-algebra-i:computing-eigenspaces}

% Sources: IMG_9323--IMG_9326; October 15 PDF, pp. 11--16.
\subsection{Diagonal and triangular matrices}

\begin{example}
  For
  \[
    D=\begin{bmatrix}a&0&0\\0&b&0\\0&0&c\end{bmatrix},
  \]
  the standard basis vectors satisfy
  \(De_1=ae_1\), \(De_2=be_2\), and \(De_3=ce_3\).
  The characteristic polynomial is
  \[
    p_D(x)=\det\begin{bmatrix}
      a-x&0&0\\0&b-x&0\\0&0&c-x
    \end{bmatrix}=(a-x)(b-x)(c-x).
  \]
  Thus the eigenvalues are the diagonal entries.
\end{example}

A triangular matrix also has its diagonal entries as eigenvalues, since
subtracting \(xI\) preserves triangularity. For instance,
\[
  T=\begin{bmatrix}a&0&0\\b&c&0\\d&e&f\end{bmatrix}
  \quad\Longrightarrow\quad
  p_T(x)=(a-x)(c-x)(f-x).
\]
Here \(Te_3=fe_3\), so \(e_3\) is an eigenvector for \(f\).
Finding eigenvectors for \(a\) and \(c\) requires solving the corresponding
kernel equations. In general, for triangular \(T=(T^i_j)\),
\[
  \det T=T^1_1\cdots T^n_n,\qquad
  p_T(x)=\prod_{j=1}^n(T^j_j-x).
\]

% !TEX root = ../../hw06.tex
% Source: supplied Homework 6 PDF, page 1, Question 1.
\begin{exercise}
  \label{ex:hw06:triangular-matrix-eigenvectors}
  Find all eigenvectors of the matrix
  \[
    M=\begin{bmatrix}
      1&1&0\\
      0&2&1\\
      0&0&2
    \end{bmatrix}.
  \]
  \solutionlink{hw06:triangular-matrix-eigenvectors}
\end{exercise}


% !TEX root = ../../hw06.tex
% Source: supplied Homework 6 PDF, page 1, Question 2.
% The source calls L a map and implicitly uses its linearity; do not add a
% hypothesis or a solution without the author's clarification.
\begin{exercise}
  \label{ex:hw06:eigenvectors-in-a-basis}
  Let \((e_1,e_2,e_3)\) be a basis of \(V\) and \(L:V\to V\) be the map
  given by
  \[
    \begin{aligned}
      L(e_1)&=e_1+e_2\\
      L(e_2)&=2e_2+e_3\\
      L(e_3)&=2e_3.
    \end{aligned}
  \]
  Find the eigenvectors of \(L\) by solving
  \[
    L(v)=\lambda v,
  \]
  where
  \[
    v=e_1a^1+e_2a^2+e_3a^3,
  \]
  directly for \(a^1,a^2,a^3,\lambda\).
  \solutionlink{hw06:eigenvectors-in-a-basis}
\end{exercise}


\subsection{A matrix that interchanges coordinates}

\begin{example}
  \label{ex:linear-algebra-i:coordinate-swap-eigenspaces}
  Let
  \[
    A=\begin{bmatrix}0&1\\1&0\end{bmatrix}.
  \]
  Its characteristic polynomial is
  \(p_A(x)=x^2-1=(x-1)(x+1)\). Row reduction gives
  \[
    A-\lambda I=
    \begin{bmatrix}-\lambda&1\\1&-\lambda\end{bmatrix}
    \sim\begin{bmatrix}1&-\lambda\\-\lambda&1\end{bmatrix}
    \sim\begin{bmatrix}1&-\lambda\\0&1-\lambda^2\end{bmatrix}.
  \]
  If \(1-\lambda^2\ne0\), the matrix reduces to the identity and its
  kernel is zero. If \(\lambda=1\), the kernel equation is
  \(v^1-v^2=0\), giving
  \[
    E_1(A)=\left\{\begin{bmatrix}t\\t\end{bmatrix}\;\middle|\;t\in\mathbb F\right\}
      =\operatorname{span}\left\{\begin{bmatrix}1\\1\end{bmatrix}\right\}.
  \]
  If \(\lambda=-1\), the equation is \(v^1+v^2=0\), giving
  \[
    E_{-1}(A)=\left\{\begin{bmatrix}t\\-t\end{bmatrix}\;\middle|\;t\in\mathbb F\right\}
      =\operatorname{span}\left\{\begin{bmatrix}1\\-1\end{bmatrix}\right\}.
  \]
  The nonzero vectors in these spaces are the eigenvectors.
  Over \(\R\), their two directions appear in
  \autoref{fig:linear-algebra-i:swap-eigendirections}.
  \begin{marginfigure}
    \centering
    % Exact eigendirections of A=[0 1;1 0] over R.
\begin{tikzpicture}[
  x=10mm,y=10mm,font=\footnotesize,line width=0.45pt,
  vector/.style={-{Latex[length=1.6mm,width=1.1mm]}},
  axis/.style={vector,black!30,line width=0.35pt},
  direction/.style={dashed,black!55,line width=0.35pt}
]
  \draw[axis] (-1.7,0) -- (1.8,0) node[right=2pt] {\(x\)};
  \draw[axis] (0,-1.6) -- (0,1.8) node[above=2pt] {\(y\)};
  \draw[direction] (-1.5,-1.5) -- (1.5,1.5) node[above right=1pt] {\(E_1\)};
  \draw[direction] (-1.5,1.5) -- (1.5,-1.5) node[below right=1pt] {\(E_{-1}\)};
  \draw[vector] (0,0) -- (1,1) node[right=3pt] {\(v=Av\)};
  \draw[vector] (0,0) -- (1,-1) node[right=3pt] {\(w\)};
  \draw[vector] (0,0) -- (-1,1) node[left=3pt] {\(Aw\)};
  \node[below left=2pt] at (0,0) {\(0\)};
  \path (current bounding box.south) ++(0,-2mm);
\end{tikzpicture}

    \caption[Eigendirections of the coordinate swap]{The map
      \(A(x,y)=(y,x)\) fixes the line \(E_1\) and reverses vectors on
      \(E_{-1}\). Here \(v=(1,1)\), \(w=(1,-1)\), and \(Aw=-w\).}
    \label{fig:linear-algebra-i:swap-eigendirections}
  \end{marginfigure}
\end{example}
% IMG_9325's final vector for -1 has a sign inconsistent with v1=-v2.
% Use the preceding kernel equation and PDF p.15.

\subsection{A matrix with no real eigenvalue}

\begin{example}
  \label{ex:linear-algebra-i:rotation-no-real-eigenvalue}
  Consider
  \[
    B=\begin{bmatrix}0&-1\\1&0\end{bmatrix},\qquad
    B\begin{bmatrix}x\\y\end{bmatrix}=\begin{bmatrix}-y\\x\end{bmatrix}.
  \]
  Its characteristic polynomial is \(p_B(x)=x^2+1\), and
  \[
    B-\lambda I=
    \begin{bmatrix}-\lambda&-1\\1&-\lambda\end{bmatrix}
    \sim\begin{bmatrix}1&-\lambda\\-\lambda&-1\end{bmatrix}
    \sim\begin{bmatrix}1&-\lambda\\0&-1-\lambda^2\end{bmatrix}.
  \]
  For real \(\lambda\), the last diagonal entry is strictly negative,
  so the kernel is zero. Thus \(B\) has no real eigenvalues.
  Over \(\C\), the roots of \(x^2+1\) are \(i\) and \(-i\).
  Geometrically, \(B\) makes a counterclockwise quarter-turn in the
  real plane, as in
  \autoref{fig:linear-algebra-i:rotation-no-real-eigendirection}.
  \begin{marginfigure}
    \centering
    % Exact action B(x,y)=(-y,x). The sample direction is schematic, not a new example.
\begin{tikzpicture}[
  x=9mm,y=9mm,font=\footnotesize,line width=0.45pt,
  vector/.style={-{Latex[length=1.6mm,width=1.1mm]}},
  axis/.style={vector,black!35,line width=0.35pt}
]
  \coordinate (O) at (0,0);
  \coordinate (V) at (1.8,0.6);
  \coordinate (BV) at (-0.6,1.8);
  \draw[axis] (-1.1,0) -- (2.4,0) node[right=2pt] {\(x\)};
  \draw[axis] (0,-0.5) -- (0,2.4) node[above=2pt] {\(y\)};
  \draw[vector] (O) -- (V) node[right=3pt] {\(v\)};
  \draw[vector] (O) -- (BV) node[above left=3pt] {\(Bv\)};
  \draw[vector,black!60] (18.43495:1.1) arc[start angle=18.43495,end angle=108.43495,radius=1.1];
  \node at (63.43495:0.7) {\(90^\circ\)};
  \node[below left=2pt] at (O) {\(0\)};
  \path (current bounding box.south) ++(0,-2mm);
\end{tikzpicture}

    \caption[A quarter-turn in the real plane]{For every real nonzero
      \(v\), the vector \(Bv\) is perpendicular to \(v\), so it cannot
      be a real scalar multiple of \(v\).}
    \label{fig:linear-algebra-i:rotation-no-real-eigendirection}
  \end{marginfigure}
\end{example}
